% Relatório do Trabalho Prático 2 — Front-end da Locadora de Veículos
% Modelo SBC: baixe sbc-template.sty e sbc-template.bst em https://www.sbc.org.br/documentos-da-sbc/
% (ou use o projeto "SBC Conference Template" do Overleaf) e coloque na mesma pasta deste arquivo.
\documentclass[12pt]{article}

\usepackage{sbc-template}
\usepackage{graphicx,url}
\usepackage[utf8]{inputenc}
\usepackage[brazil]{babel}
\usepackage[T1]{fontenc}
\usepackage{booktabs,array}

\sloppy

% Mostra a figura se o arquivo existir; senão, uma caixa avisando para inserir o print.
\newcommand{\figura}[3]{%
  \begin{figure}[ht]\centering
  \IfFileExists{figuras/#1}{\includegraphics[width=\textwidth]{figuras/#1}}%
  {\fbox{\parbox{0.9\textwidth}{\centering\vspace{1.5cm}\textbf{[inserir print: figuras/#1]}\vspace{1.5cm}}}}
  \caption{#2}\label{#3}
  \end{figure}}

\title{Locadora de Veículos: front-end web integrado à API}

\author{Paula Bessa\inst{1}}

\address{Pontifícia Universidade Católica de Minas Gerais (PUC Minas)\\
  Belo Horizonte -- MG -- Brasil
  \email{paulabessa566@gmail.com}
}

\begin{document}

\maketitle

\section{O problema}

Uma locadora de veículos precisa controlar a frota, os clientes e os aluguéis: quais carros estão
disponíveis, quem está com qual veículo, quando ele deve voltar e quanto cada cliente já gerou de
faturamento. No Trabalho Prático~1 foi construída a API que guarda e consulta esses dados
(C\#, ASP.NET Core, Entity Framework e SQL Server Express). Faltava uma interface para que um
atendente use o sistema sem precisar do Swagger. Este trabalho entrega esse front-end, que consome
as rotas reais da API, sem dados fixos no código.

\section{Modelagem adotada}

O banco tem cinco entidades: \textbf{Fabricante}, \textbf{Categoria}, \textbf{Veículo},
\textbf{Cliente} e \textbf{Aluguel}. Um fabricante e uma categoria têm vários veículos; um cliente e
um veículo têm vários aluguéis. A categoria guarda a diária base. O aluguel guarda as datas de
retirada, devolução prevista e devolução, a quilometragem inicial e final, o valor da diária e o
valor total, calculado na devolução. O veículo tem situação: Disponível, Alugado ou Manutenção.
O código e o relatório da API estão no repositório do Trabalho Prático~1 (\url{https://github.com/paulabessa21/Trabalho-Pratico-Tecnologias-para-ADS}).

\section{Arquitetura do front-end}

O front-end fica na pasta \texttt{Frontend/} do mesmo repositório. É feito com HTML, CSS (Bootstrap~5)
e JavaScript puro, sem etapa de build. A comunicação com a API usa somente \texttt{fetch}, em um único
arquivo (\texttt{js/api.js}) que também traduz os erros da API (400, 404, 409 e 500) em mensagens
para o usuário. As telas são descritas em \texttt{js/entidades.js} e desenhadas por um motor comum
(\texttt{js/crud.js}), o que mantém o mesmo comportamento de lista, filtro, formulário e exclusão
em todas as entidades. Para o navegador poder chamar a API em outro endereço, foi acrescentada à
\texttt{Program.cs} uma política de CORS.

\section{Wireframes}

Os wireframes foram feitos em média fidelidade, em tons de cinza, antes da implementação. As Figuras
\ref{fig:w1} a \ref{fig:w6} mostram as seis telas. Link do projeto no Figma: \textbf{[colar o link compartilhável aqui]}.

\figura{wireframe-01-painel.png}{Painel com resumo e duas consultas da API.}{fig:w1}
\figura{wireframe-02-lista-veiculos.png}{Lista de veículos com filtros.}{fig:w2}
\figura{wireframe-03-modal-veiculo.png}{Formulário de cadastro e edição em modal.}{fig:w3}
\figura{wireframe-04-clientes.png}{Lista de clientes com filtros.}{fig:w4}
\figura{wireframe-05-alugueis.png}{Lista de aluguéis, filtros e registro de devolução.}{fig:w5}
\figura{wireframe-06-devolucao-exclusao.png}{Modais de devolução e de confirmação de exclusão.}{fig:w6}

\section{Funcionalidades e rotas utilizadas}

A Tabela~\ref{tab:rotas} relaciona cada tela às rotas da API e aos filtros oferecidos. Todas as
entidades têm cadastro, edição (exceto aluguéis, que na API só têm criação, devolução e exclusão),
exclusão com confirmação e pesquisa com filtros em pelo menos dois campos.

\begin{table}[ht]
\centering\small
\footnotesize
\begin{tabular}{p{1.7cm}p{4.7cm}p{5.0cm}}
\toprule
Tela & Rotas & Filtros \\
\midrule
Veículos & GET/POST \path{/api/veiculos}; PUT/DELETE \path{/{id}}; GET \path{/disponiveis}; GET \path{/com-total-alugueis} & modelo ou placa, situação, categoria, fabricante, ano; ``só disponíveis'' usa a rota \path{/disponiveis} \\
Clientes & GET/POST \path{/api/clientes}; PUT/DELETE \path{/{id}}; GET \path{/sem-aluguel} & nome, CPF, e-mail; ``sem aluguel'' usa a rota \path{/sem-aluguel} \\
Aluguéis & GET/POST \path{/api/alugueis}; PATCH \path{/{id}/devolucao}; DELETE \path{/{id}}; GET \path{/em-aberto} & texto, cliente, situação, período de retirada; ``em aberto'' usa \path{/em-aberto?clienteId=} \\
Fabricantes & GET/POST \path{/api/fabricantes}; PUT/DELETE \path{/{id}} & nome, código \\
Categorias & GET/POST \path{/api/categorias}; PUT/DELETE \path{/{id}} & nome, diária mínima e máxima \\
Painel & GET \path{/api/alugueis/faturamento-por-periodo}; GET \path{/api/veiculos/com-total-alugueis}; GET \path{/api/alugueis/em-aberto} & período do faturamento \\
\bottomrule
\end{tabular}
\caption{Telas, rotas e filtros.}\label{tab:rotas}
\end{table}

\subsection{Duas consultas, do banco ao front-end}

\textbf{Consulta 1: faturamento por cliente.} A rota \texttt{GET /api/alugueis/faturamento-por-periodo}
faz um \texttt{INNER JOIN} entre aluguéis e clientes, agrupa por cliente e soma o valor dos aluguéis
devolvidos no período. O painel envia \texttt{dataInicio} e \texttt{dataFim} e mostra o resultado em
tabela, com o total. \textbf{[inserir print com os dados reais: figuras/tela-painel.png]}

\textbf{Consulta 2: veículos com total de aluguéis.} A rota \texttt{GET /api/veiculos/com-total-alugueis}
usa um \texttt{LEFT JOIN} com agrupamento, e por isso lista também os veículos que nunca foram alugados
(total zero). O painel mostra cada veículo com uma barra proporcional ao total.

\figura{tela-painel.png}{Painel da aplicação com dados reais.}{fig:painel}
\figura{tela-veiculos.png}{Lista de veículos com filtro ``só disponíveis''.}{fig:veiculos}
\figura{tela-alugueis.png}{Aluguéis em aberto e registro de devolução.}{fig:alugueis}

\section{Regras de negócio refletidas na interface}

\begin{itemize}
  \item Ao criar um aluguel, a lista de veículos mostra só os disponíveis; a quilometragem inicial e a
        diária são sugeridas a partir do veículo escolhido e da sua categoria.
  \item A devolução só aceita quilometragem final maior ou igual à inicial. A API calcula o valor total e
        libera o veículo.
  \item Excluir fabricante, categoria, cliente ou veículo em uso retorna erro 409 da API, e a mensagem
        aparece na própria janela de confirmação.
  \item Campos têm validação no navegador (CPF com 11 dígitos, placa no padrão antigo ou Mercosul, e-mail)
        e a API continua validando.
\end{itemize}

\section{Como executar}

Com o SQL Server Express ativo, execute \texttt{dotnet run} na pasta do backend (a API sobe em
\texttt{http://localhost:5202}) e abra \texttt{Frontend/index.html} no navegador. Se a porta for outra,
altere \texttt{API\_URL} em \texttt{Frontend/js/config.js}.

\section{Conclusão}

O front-end cobre todas as funcionalidades projetadas no Trabalho Prático~1 e usa as rotas da Etapa~3.
O ponto que mais exigiu cuidado foi traduzir os erros da API (validação, conflito e falha de servidor)
em mensagens claras para quem usa o sistema.

\end{document}
